\documentclass[10pt]{article}
\usepackage[utf8]{inputenc}
\usepackage[a4paper,height=20cm,width=13cm]{geometry}
\usepackage{amssymb}
\usepackage{dsfont}
\usepackage{calc}
\usepackage{graphicx}
% \usepackage{pst-node}
\usepackage{fourier}
\usepackage{euscript}
\usepackage{amsmath,amsthm}
\def\phi{\varphi}
\def\theta{\vartheta}
\def\P{\EuScript P}
\def\M{\EuScript M}
\def\N{\EuScript N}
\def\B{\EuScript B}
\def\C{\EuScript C}
\def\D{\EuScript D}
\def\U{\EuScript U}
\def\V{\EuScript V}
\def\S{\EuScript S}
\def\sm{\smallsetminus}
\def\niff{\nleftrightarrow}
\def\ZZ{\mathds Z}
\def\NN{\mathds N}
\def\PP{\mathds P}
\def\QQ{\mathds Q}
\def\RR{\mathds R}
\def\<{\langle}
\def\>{\rangle}
\def\E{\exists}
\def\A{\forall}
\def\0{\varnothing}
\def\imp{\rightarrow}
\def\iff{\leftrightarrow}
\def\IMP{\Rightarrow}
\def\IFF{\Leftrightarrow}
\def\range{\textrm{im}}
\def\Mod{\textrm{Mod}}
\def\Aut{\textrm{Aut}}
\def\Th{\textrm{Th}}
\def\acl{\textrm{acl}}
\def\dcl{\textrm{dcl}}
\def\dom{{\rm dom}}
\def\eq{{\rm eq}}
\def\fp{{\rm fp}}
\def\tp{\textrm{tp}}
\def\equivL{\stackrel{\smash{\scalebox{.5}{\rm L}}}{\equiv}}


\def\nonfork{\mathop{\raise0.2ex\hbox{\ooalign{\hidewidth$\vert$\hidewidth\cr\raise-0.9ex\hbox{$\smile$}}}}}

\newcommand{\labella}[1]{{\sf\footnotesize #1}\hfill}
\renewenvironment{itemize}
  {\begin{list}{$\triangleright$}{%
   \setlength{\parskip}{0mm}
   \setlength{\topsep}{0mm}
   \setlength{\rightmargin}{0mm}
   \setlength{\listparindent}{0mm}
   \setlength{\itemindent}{0mm}
   \setlength{\labelwidth}{3ex}
   \setlength{\itemsep}{0mm}
   \setlength{\parsep}{0mm}
   \setlength{\partopsep}{0mm}
   \setlength{\labelsep}{1ex}
   \setlength{\leftmargin}{\labelwidth+\labelsep}
   \let\makelabel\labella}}{%
   \vspace*{-.3\baselineskip}
  \end{list}}
%\def\ssf#1{\textsf{\small #1}}
\newcounter{ex}
\newenvironment{exercise}{\bigskip\addtocounter{ex}{1}\textbf{Esercizio \theex.\ }}{}
\pagestyle{empty}
\parindent0ex
\parskip1ex
\raggedbottom
\def\nsR{{}^*\!\RR}
\def\ssf#1{\textsf{#1}}
\renewcommand{\baselinestretch}{1.3}


\usepackage{fancyhdr}
\pagestyle{fancy}
\lfoot{Teoria dei Modelli a.a.~2026/27}
\rhead{}
\rfoot{{\small\thepage}}
\cfoot{}

%
% INDICARE NOME E COGNOME DI TUTTI GLI AUTORI
%
\lhead{Nome/i Cognome/i}
\begin{document}

\begin{exercise}(Beatrice Calderaro)
  Let $M$ be a semigroup $M$ with the property that $M\cdot a$ is infinite for every $a\in M$.
  Prove that for every finite coloring of $M$ there is a sequence $\bar a=\langle a_i:i<\omega\rangle$ with an infinite range such that $\fp\,\bar a$ is monochromatic.
  Hint: prove that there is an idempotent not in $M$. Fix a minimal left ideal of $\M$. If all idempotents of $\M$ were in $M$ then $\M\sm M$ would be an idempotent: a contradiction.
  (You need not rewrite the part of the proof that are in the notes.)
\end{exercise}

\begin{exercise}(Martina Apicella)
  Let $M$ be a semigroup $M$ with the property that $a\cdot M$ is infinite for every $a\in M$.
  Prove that for every finite coloring of $M$ there is a sequence $\bar a=\langle a_i:i<\omega\rangle$ with an infinite range such that $\fp\,\bar a$ is monochromatic.
  Hint: similar to the previous exercise but proving that $\M\not\nsubseteq M$ requires a longer argument. Recall that $\M*b$ is a minimal left ideal for every $b$.
\end{exercise}

\begin{exercise}(Mattia Lo Mele \& Andrea Morello)
  Let $\M$ be a type-definable minimal left ideal.
  Let $u\in\M$ be an idempotent.
  Prove that $u*\U$ is a minimal right ideal.
\end{exercise}

\begin{exercise}(Filippo Giugie \& Daniele Lanciotti)
Prove that 

 \hfil$\displaystyle\bigcup\big\{\M\ :\ \M\text{ a minimal left ideal}\big\} =\bigcup\big\{\M\ :\ \M\text{ a minimal right ideal}\big\}$.

\end{exercise}

\begin{exercise}(Francesco Dam\&	Gabriele Maiolino)
  Let $\M$ and $\N$ be type-definable minimal left ideals.
  Let $u\in\M$ and $v\in\N$ be idempotents.
  Prove that $u*\M$ and $v*\N$ are isomorphic groups.
  Hint: the isomorphism is given by the map $a\mapsto v*a$ with inverse $a\mapsto u*a$ works well if we can assume that $u$ and $v$ are such that $u*v=u$ and $v*u=v$.
\end{exercise}

\end{document}


